% Chapter 7 -- translated by Claude (2026), continuing the voice and
% register of the Burkett-Hutcheson chapters.

\chapter{Two-ports}

\section{Introduction}

A two-port is a network that has two pairs of terminals. In
Symbulator, two-ports can be simulated as circuit elements, assuming
that:

\begin{itemize}
\item the network is composed only of linear elements,
\item the network contains no independent sources (it may contain
  dependent sources),
\item the two upper nodes of the terminals are connected to distinct
  nodes of the circuit,
\item the two lower nodes of the terminals are considered connected
  to the reference node, 0.
\end{itemize}

Symbulator accepts two-ports with parameters of 6 different types:

\begin{itemize}
\item admittance parameters
\item impedance parameters
\item hybrid parameters
\item gain (inverse hybrid) parameters
\item transmission parameters
\item inverse transmission parameters
\end{itemize}

There are three procedures related to two-ports that can be done in
Symbulator:

\begin{enumerate}
\item Simulate circuits that contain two-ports. For this, two-ports
  are used as circuit elements.
\item Find the gains in a two-port. For this, the \code{gain} tool
  is used after a circuit simulation.
\item Find the equivalent two-port of a network. For this, the
  \code{port} tool is used with the description of a network.
\end{enumerate}

In this chapter we will learn everything related to two-ports in
Symbulator.

\section{Notation}

In the case of two-ports, the minimum information necessary to
describe them completely is the following:

\emph{Identification.} The letters that identify the different
two-ports are:

\begin{itemize}
\item impedance parameters, the letter \code{z}
\item admittance parameters, the letter \code{y}
\item hybrid parameters, the letter \code{h}
\item gain parameters, the letter \code{g}
\item transmission parameters, the letter \code{a}
\item inverse transmission parameters, the letter \code{b}
\end{itemize}

\emph{Name.} Any name is valid that begins with the letter
identifying the two-port you want to describe, and that is not a
reserved variable, like those listed in the User's Manual and in
section 1.9.4 of this book. Example: \code{zc} and \code{yc} are
reserved variables that cannot be used as names.

\emph{Nodes.} A two-port has four connection points, two upper and
two lower. For a simulation in Symbulator, the two lower ones are
always assumed connected to the reference node. The upper ones must
be connected to nodes distinct from each other. It is only necessary
to give the simulator the names of these two upper nodes to which
the two-port is connected.

So, two-ports are defined as shown below:

\begin{entry}
lettername, node 1, node 2
\end{entry}

As necessary, one or two zeros are added at the end.

\emph{Value of the parameters.} Two-ports are unique elements,
because they require something no other element requires in
Symbulator: variables stored previous to the simulation. Each
two-port that is simulated as part of a circuit requires that the
values of its 4 parameters be stored in 4 variables, one for each
parameter, in the following form:

\begin{itemize}
\item parameter 11, in the variable \code{lettername11}
\item parameter 12, in the variable \code{lettername12}
\item parameter 21, in the variable \code{lettername21}
\item parameter 22, in the variable \code{lettername22}
\end{itemize}

Take care to enter these values in whole units, and not in their
multiples. This means they must not be entered in milli (m), micro
(µ), kilo (k), mega (M), etc.

\section{Related answers}

In the case of a two-port, two related answers are given. They are
two currents: the current entering the two-port through the first
node, and the current entering through the second node. Remember
that these two nodes are the upper ones, since the lower one on both
sides is the reference node. These currents are stored in variables
called \code{inameNODE}, that is, \code{i} followed by the name of
the two-port and the name of the node. So, for a two-port called
\code{yp}, connected to nodes 1 and 2, the related answers will be
\code{iyp1} and \code{iyp2}. Remember that both currents are
considered as entering the two-port through the upper nodes.

Let us see our first two-port problem.

\problem{039}

\emph{Problem: In the following circuit, the parameters of the
two-port shown are $z_{11} = 20/3$, $z_{12} = 1/3$,
$z_{21} = 2500/3$, $z_{22} = 200/3$. Find: a) the voltage at the
load node and the currents entering the two-port at both ends, b)
the input impedance and the current, voltage and power gains of the
two-port, c) the output impedance of the network, seen by the load,
and d) the equivalents of this two-port in the other five types of
parameters.}

Solution:

This problem, specially designed to inaugurate our experience with
two-ports, will be solved in four parts.

Question a) will be answered by simulating the circuit using a
direct current analysis. We will call the two-port \code{zp}. First,
we define the variables that describe the parameters of the
two-port:

\begin{entry}
20/3→zp11:1/3→zp12:2500/3→zp21:200/3→zp22
\end{entry}

After entering this line and pressing \key{ENTER}, we will see in
the history area the value of the last of these parameters as the
answer to our command. What interests us is that the values were
stored in the variables. Note that we used the separator \code{:}
and the \key{STO} key. Once the two-port parameter variables are
stored, we proceed to order the simulation. Below, the description
we use:

\begin{entry}
sq\dc("e1,1,0,1;r1,1,2,2;r2,3,0,20;zp,2,3,0")
\end{entry}

Note that, since the description of the two-port only requires three
terms, we fill with a zero to keep the symmetry. We press
\key{ENTER}. After the status messages, `Done' appears. The problem
has asked us for the voltage at the load node and the currents
entering the two-port at both ends. We obtain these values like
this: the voltage at the load node we obtain with \code{v3}, and it
is worth 2500/71 volts; the current entering the two-port on the
left we obtain with \code{izp2}, and it is worth 13/71 amperes; and
the current entering the two-port on the right we obtain with
\code{izp3}, and it is worth $-125/71$ amperes. The negative sign of
this last current tells us that the current is really leaving the
two-port on the right, not entering.

We have learned that a two-port can be easily simulated in
Symbulator as one more circuit element. To answer question b), we
will use a new tool.

\section{The gain tool}

This tool helps us find the current, voltage and power gains, and
the input impedance of a network. This network may or may not
contain two-ports. Let us learn to use this tool at the same time we
answer the second question. The \code{gain} tool is executed like
this:

\begin{entry}
sq\gain()
\end{entry}

It takes no input arguments between the parentheses. However, when
executing it, a form opens in which we must enter four input values:

\emph{The voltage at the input terminal node.} In the case of our
problem, this voltage is \code{v2}, because this is the voltage at
the input node of the two-port, the one toward the source.

\emph{The current entering through the input terminal.} In the case
of our problem, this current is \code{izp2}, because this is the
current entering the two-port on the input side, that is, the side
toward the source.

\emph{The voltage at the output terminal node.} In the case of our
problem, this voltage is \code{v3}, because this is the voltage at
the output node of the two-port, the one toward the load.

\emph{The current entering through the output terminal.} In the case
of our problem, this current is \code{izp3}, because this is the
current entering the two-port on the output side, that is, the side
toward the load. Note that saying output current does not mean a
current that is leaving. The current must be defined as entering the
two-port, but on the output side.

We press \key{ENTER}, twice if necessary. Next a screen appears,
which works in a way similar to that of the \code{thevenin} tool, in
which we must press \key{ENTER} as many times as there are values we
seek.

We see on the screen that the voltage gain, called $A_v$ or $G_v$
depending on the textbook, is worth 500/9; the current gain, called
$A_i$ or $G_i$, is worth $-125/13$; the power gain, called $A_p$ or
$G_p$, is worth 62500/117; and the input impedance, called $Z_i$ or
$R_i$, is worth 45/13. Besides showing these values on screen, the
tool also stores them in variables with the names \code{av},
\code{ai}, \code{ap}, and \code{zi}, respectively. All these values,
being stored in memory, can be requested later, even in the form of
approximate values if desired.

We have learned the use of the \code{gain} tool. Let us now answer
the next question.

\section{Output impedance}

Question c) requires some theoretical knowledge. Theory teaches us
that the output impedance of the network, seen by the load, is the
equivalent impedance of the network, from the terminals where the
load is connected, when the load has been removed and the
independent sources have been killed. To kill an independent voltage
source is to turn it into a short circuit. To kill an independent
current source is to turn it into an open circuit. We already
defined the parameters of this two-port when we answered the first
question, and they are in memory. Therefore, there is no need to
define them again. Below, the description we use:

\begin{entry}
sq\thevenin("e1,1,0,0;r1,1,2,2;zp,2,3,0",3,0)
\end{entry}

Note that the load has been removed from the description, and that
the source has been killed, that is, its value has been made null.
This is the same as having turned it into a short circuit. We press
\key{ENTER}, choose Passive and DC, and press \key{ENTER}. On the
answer screen, it appears that \code{zth} is 450/13 ohms, which is
the correct output impedance.

We have learned how to find the output impedance of a network. We
have also seen that the \code{thevenin} tool works with two-ports
too. Let us now answer the last question, using a new tool.

\section{The port tool}

This tool gives us the parameters of the equivalent two-port of any
network with two terminals. Let us learn to use this tool, at the
same time we answer the fourth question. The \code{port} tool is
executed like this:

\begin{entry}
sq\port(description of the network,node1,node2)
\end{entry}

Three input arguments are given between the parentheses:

\emph{The description of the two-terminal network} whose two-port
equivalent you want to find. This network must be a circuit composed
only of passive elements, and it cannot contain independent sources.
This network may be composed, say, of resistances only. It may also
be composed of one or more two-ports. Or of a mixture of resistive
elements and two-ports. Since this network can contain other
two-ports, the \code{port} tool is a powerful ally for finding the
equivalent of a given two-port in other, distinct parameters, and
for reducing a set of several two-ports (in series, in parallel or
in combination) to a single equivalent two-port.

\emph{The upper node of the input terminal of the network.}

\emph{The upper node of the output terminal of the network.}

It is assumed that the lower node of the two terminals is the
reference node.

Once we have written the line shown, we press \key{ENTER}. When it
executes, a form opens in which we must make some selections. The
three selections are:

\emph{The type of parameter we want to find.} The options are
\code{z}, \code{y}, \code{h}, \code{g}, \code{a} and \code{b}, whose
meanings we already listed in section 7.2.

\emph{The name we wish to give the two-port,} without including the
letter that identifies the type of two-port. For example, if we want
the two-port to be called \code{zp1}, the name we must enter here is
only \code{p1}, since the \code{z} we have already chosen above. We
can use any name, except those that, when combined with the two-port
letter we have chosen, turn out to be the name of a reserved
variable. For example, if we choose a two-port of type \code{z} or
\code{y}, we cannot use a name like \code{c}, because the complete
name of the two-port would be \code{zc} or \code{yc}, which are both
reserved variables, as we pointed out in section 1.9.4. There is a
large number of reserved variables that begin with \code{z} and with
\code{y}. A complete list can be found in the User's Manual. As a
general recommendation, use a letter or a number as the name, and
let that letter not be the letter \code{c}.

\emph{The type of analysis we want.} We can choose between direct
current analysis DC, alternating current AC and frequency domain
FD.

Next, a simulation of the network is executed. Then an answer screen
appears that works in a way similar to the previous ones. We must
press \key{ENTER} four times to see the four parameters sought,
presented on the screen.

These parameters will also be stored in the memory of the
calculator, in variables with the names that correspond to the name
of the two-port, according to the naming convention we explained in
section 7.2.

With this knowledge, let us now see how we can answer question d)
using the \code{port} tool. The question asks us to find the
equivalents of the two-port in all the other types of parameters.
Below, the solution to this question.

Let us first find the equivalent in admittance parameters of the
given two-port. Enter the following command line in the entry line:

\begin{entry}
sq\port("zp,1,2,0",1,2)
\end{entry}

This line orders the use of the \code{port} tool to find the
equivalent of the network described, which is the two-port
\code{zp}, whose impedance parameters are already known and are
called \code{zp11}, \code{zp12}, \code{zp21} and \code{zp22}. These
parameters are already stored in memory, from when we answered the
first question. If you have any doubt about whether these parameters
are stored in memory or not, you can use the Var-Link screen to
verify. If these parameters were deleted, it is necessary to enter
them again, as shown in the earlier question. The fact is that the
parameters of a two-port must be stored before executing any
simulation involving that two-port.

The two-port is connected between nodes 1 and 2, and this is
indicated in the command line. Although it may seem unnecessary in
this case, the two end nodes of the network to be reduced must be
specified, because this network will not always be composed of a
single element, as we will see later.

So, we press \key{ENTER}. We specify \code{y} as the type of
two-port desired, \code{p} as the name and DC as the analysis. The
answers are shown on the screen: the four admittance parameters.
This is the correct equivalent. Remember that these parameters, in
addition to appearing on screen, are also stored in variables with
the corresponding names. Let us now find another equivalent of the
two-port in another type of parameters.

To find the equivalent in hybrid parameters of the given two-port,
the procedure is identical. Just for didactic purposes, we will now
use as the starting two-port the recently found two-port \code{yp}
whose parameters are already stored in memory, instead of the
two-port \code{zp} we used in the previous line. Enter the following
command line:

\begin{entry}
sq\port("yp,1,2,0",1,2)
\end{entry}

We press \key{ENTER}. We specify \code{h} as the type of two-port
desired, \code{p} as the name and DC as the analysis. The answers
are shown.

To find the equivalent in gain parameters of the given two-port, the
procedure is the same. We will use as the starting two-port the
two-port \code{hp}, whose parameters are in memory. Enter the
following command line:

\begin{entry}
sq\port("hp,1,2,0",1,2)
\end{entry}

We press \key{ENTER}. We specify \code{g} as the type of two-port
desired, \code{p} as the name and DC as the analysis. The answers
are shown.

Let us find the equivalent in transmission parameters of the
two-port, using as the starting two-port the two-port \code{gp}.
Enter the following command line:

\begin{entry}
sq\port("gp,1,2,0",1,2)
\end{entry}

We press \key{ENTER}. We specify \code{a} as the type of two-port
desired, \code{p} as the name and DC as the analysis. The answers
are shown.

Let us find the equivalent in inverse transmission parameters of the
two-port, using as the starting two-port the two-port \code{ap}.
Here we find a variant in the procedure. When we solved question b),
a moment ago, the \code{gain} tool generated a variable called
\code{ap}, which contained the power gain. This variable is still in
memory. If we forget this, and try to execute the tool using
\code{ap} as the name of the two-port, we will receive an error from
Symbulator telling us that the name we have assigned to element \#1
of our circuit description is not appropriate, and that we should
please correct this description. The name is inappropriate because
it is the name of a variable in use.

Therefore, we must delete the variable before executing the tool. We
can put both commands in a single line, using the separator. For
example, enter the following command line:

\begin{entry}
DelVar ap:sq\port("ap,1,2,0",1,2)
\end{entry}

Since the variable \code{ap} has been deleted, now we can use this
variable as the name of the two-port. We press \key{ENTER}. We
specify \code{b} as the type of two-port desired, \code{p} as the
name and DC as the analysis. The answers are shown.

So, we have found the five equivalents of the two-port that was
given to us.

\section{Verification}

But how could we verify that these two-ports are the correct
equivalents of our original two-port? Simple. Let us find the
equivalent in impedance parameters, that is, in the original type,
using as the starting two-port the two-port \code{bp}, which was the
last equivalent found. Any error generated along the way will impact
our final answer. So, if we now obtain the correct equivalent, it
will mean that all the intermediate equivalents were correct. Enter
the following command line:

\begin{entry}
sq\port("bp,1,2,0",1,2)
\end{entry}

We press \key{ENTER}. We specify \code{z} as the type of two-port
desired, \code{p} as the name and DC as the analysis. The answers
are shown, and we can see that these impedance parameters coincide
perfectly with those the problem gave us as initial data. This means
that all the equivalents we have found are correct. Let us see more
solved problems, as practice to master the use of two-ports in
Symbulator.

\problem{040}

\emph{Problem: The parameters of the two-port shown are
$z_{11} = 20$, $z_{12} = 3$, $z_{21} = 100$, $z_{22} = 5$. Find: a)
the input impedance and the current, voltage and power gains of the
two-port, b) the output impedance of the network, seen by the
load.}

Solution: Below, the description we use for the first question:

\begin{entry}
20.→zp11:3.→zp12:100.→zp21:5.→zp22:
sq\dc("j1,0,1,is;r1,1,0,50.;zp,1,2,0;r2,2,0,100."):sq\gain()
\end{entry}

In this line are the definition of the parameters, the command to
run the simulation and the command to run the \code{gain} tool. We
press \key{ENTER}. After the status messages, the tool's form
appears, in which we enter \code{v1}, \code{izp1}, \code{v2},
\code{izp2}. We press \key{ENTER}. We receive on the screen, and in
memory, the answers: \code{ai}~$= -.952$, \code{av}~$= 5.56$,
\code{ap}~$= 5.29$, \code{zi}~$= 17.1$.

Below, the description for the second question:

\begin{entry}
sq\thevenin("r1,1,0,50.;zp,1,2,0",2,0)
\end{entry}

Note that the current source was removed. This way of killing it is
equivalent to having entered it with a null value, and has the
advantage of saving the simulator trouble. We press \key{ENTER}. We
select Passive and DC, and press \key{ENTER}. We obtain:
\code{zth}~$= .714$. These are the correct answers.

\problem{041}

\emph{Problem: Find the equivalent of the following network, in
admittance parameters.}

Solution: Below, the description we use:

\begin{entry}
sq\port("r1,1,2,10.;r2,1,0,5.;r3,2,0,20.",1,2)
\end{entry}

We press \key{ENTER}. We select \code{y}, some name and DC, and
press \key{ENTER}. We obtain the following admittance parameters:
.3, $-.1$, $-.1$, .15.

\problem{042}

\emph{Problem: Find the equivalent of the following network, in
admittance parameters.}

Solution: Below, the description we use:

\begin{entry}
sq\port("r1,1,2,10.;r2,2,3,5;r3,1,3,20;r4,2,0,40",1,3)
\end{entry}

We press \key{ENTER}. We select \code{y}, some name and DC, and
press \key{ENTER}. We obtain the following admittance parameters:
.1192, $-.1115$, $-.1115$, .1269.

\problem{043}

\emph{Problem: Find the equivalent of the following network, in
admittance parameters.}

Solution: Below, the description we use:

\begin{entry}
sq\port("s1,x,1,0;j1,1,0,.2*v2;r1,1,0,10.;j2,1,0,.5*is1;
        r2,2,1,5.",x,2)
\end{entry}

Note that we introduced a short circuit with the purpose of using
its current as the control current for source \code{j2}. We press
\key{ENTER}. We select \code{y}, some name and DC, and press
\key{ENTER}. We obtain the following admittance parameters: .6, 0.,
$-.2$, .2.

\problem{044}

\emph{Problem: The following circuit is the linear equivalent of a
transistor in the common-emitter configuration with resistive
feedback between the collector and the base. Find the equivalent, in
admittance parameters.}

Solution: Below, the description we use:

\begin{entry}
sq\port("r1,1,0,5E2;r2,1,2,2E3;j1,2,0,.0395*v1;r3,2,0,1E4",1,2)
\end{entry}

We press \key{ENTER}. We select \code{y}, some name and DC, and
press \key{ENTER}. We obtain the following admittance parameters:
.0025, $-5$.E$-4$, .039, 6.E$-4$.

\problem{045}

\emph{Problem: Find the equivalent of the following network, in
impedance parameters.}

Solution: Below, the description we use:

\begin{entry}
sq\port("r1,1,3,20.;r2,3,2,50.;r3,3,0,25.",1,2)
\end{entry}

We press \key{ENTER}. We select \code{z}, some name and DC, and
press \key{ENTER}. We obtain the following impedance parameters:
45., 25., 25., 75.

\problem{046}

\emph{Problem: Find the equivalent of the following network, in
impedance parameters.}

Solution: Below, the description we use:

\begin{entry}
sq\port("r1,1,0,40.;r2,1,3,20.;r3,3,0,25.;r4,3,2,50.",1,2)
\end{entry}

We press \key{ENTER}. We select \code{z}, some name and DC, and
press \key{ENTER}. We obtain the following impedance parameters:
21.2, 11.8, 11.8, 67.6.

\problem{047}

\emph{Problem: Find the equivalent of the following network, in
impedance parameters.}

Solution: Below, the description we use:

\begin{entry}
sq\port("r1,1,3,20.;r2,3,2,50.;r3,3,4,25.;e1,4,0,.5*v2",1,2)
\end{entry}

We press \key{ENTER}. We select \code{z}, some name and DC, and
press \key{ENTER}. We obtain the following impedance parameters:
70., 100., 50., 150.

\problem{048}

\emph{Problem: This problem has no figure. Find the input impedance
and the current, voltage and power gains of a two-port, with
impedance parameters $z_{11} = 4$, $z_{12} = 1.5$, $z_{21} = 10$ and
$z_{22} = 3$, if it has an input source $V_s$ at its input in series
with a 5~$\Omega$ resistance, while the load is 2~$\Omega$.}

Solution: Below, the command we use:

\begin{entry}
4→zp11:1.5→zp12:10→zp21:3→zp22:
sq\dc("e1,1,0,vs;r1,1,2,5;r2,3,0,2;zp,2,3,0"):sq\gain()
\end{entry}

In this line are the definition of the parameters, the command to
run the simulation and the command to run the \code{gain} tool. We
press \key{ENTER}. After the status messages, the tool's form
appears, in which we enter \code{v2}, \code{izp2}, \code{v3},
\code{izp3}. We press \key{ENTER}. We receive on the screen and in
memory the answers: \code{ai}~$= -2$., \code{av}~$= 4$.,
\code{ap}~$= 8$., \code{zi}~$= 1$.

\problem{049}

\emph{Problem: Find the equivalent of the following network, in
hybrid parameters.}

Solution: Below, the description we use:

\begin{entry}
sq\port("r1,1,3,1.;r2,3,2,6.;r3,3,0,4.",1,2)
\end{entry}

We press \key{ENTER}. We select \code{h}, some name and DC, and
press \key{ENTER}. We obtain the following hybrid parameters: 3.4,
.4, $-.4$, .1.

\problem{050}

\emph{Problem: The following circuit is the equivalent of a
high-frequency transistor. Find the equivalent, in impedance
parameters.}

Solution: Below, the description we use:

\begin{entry}
sq\port("r1,1,0,1E5;ca,1,0,5E-12;cb,1,2,1E-12;
        j1,2,0,.01*v1;r2,2,0,1E4",1,2)
\end{entry}

We press \key{ENTER}. We select \code{z}, some name and AC, and
press \key{ENTER}. As the frequency, we enter \code{1E8}, and press
\key{ENTER}. We obtain the following impedance parameters:
$89.7 - 98.4i$, $94.1 - 4.34i$, $528. + 9401.i$, $564. - 35.5i$.
Since these values are stored in memory, we can use the
\code{absang} tool to see them in polar form: 133.1∠$-$47.6°,
94.2∠$-$2.64°, 9416.∠86.8°, 565.∠$-$3.60°.

\problem{051}

\emph{Problem: Find the equivalent of the following network, in
impedance parameters.}

Solution: Below, the description we use:

\begin{entry}
sq\port("r1,1,3,8.;e1,3,0,.1*v2;j1,2,0,.05*v1;r2,2,0,12.",1,2)
\end{entry}

We press \key{ENTER}. We select \code{z}, some name and DC, and
press \key{ENTER}. We obtain the following impedance parameters:
7.55, 1.13, $-4.53$, 11.3.

\problem{052}

\emph{Problem: The following circuit is the equivalent of a
high-frequency transistor. Find the equivalent, in hybrid
parameters.}

Solution: Below, the description we use:

\begin{entry}
sq\port("r1,1,0,2E3;ca,1,0,5E-12;cb,1,2,2.5E-12;
        j2,2,0,.005*v1",1,2)
\end{entry}

We press \key{ENTER}. We select \code{h}, some name and AC, and
press \key{ENTER}. As the frequency, we enter \code{1E8}, and press
\key{ENTER}. We obtain the following hybrid parameters:
$615. - 923.i$, $.231 + .154i$, $2.85 - 4.77i$,
$.00119 + 9.62$E$-4i$. Since these values are stored in memory, we
can use the \code{absang} tool to see them in polar form:
1109∠$-$56.3°, .277∠33.7°, 5.55∠$-$59.2°, .00153∠38.9°.

\problem{053}

\emph{Problem: Find a) the equivalent of the following network, in
hybrid parameters, and b) the output impedance of the network, seen
from the perspective of the load, if the input contains $V_s$ in
series with $R_s = 200~\Omega$.}

Solution:

Below, the description we use for the first part:

\begin{entry}
sq\port("r1,1,0,1E3;j1,0,1,1E-5*v2;e1,0,3,100*v1;
        r2,3,2,1E4",1,2)
\end{entry}

We press \key{ENTER}. We select \code{h}, some name and DC, and
press \key{ENTER}. We obtain the following hybrid parameters:
1000., .01, 10., 2.E$-4$.

Below, the description we use for the second part:

\begin{entry}
sq\thevenin("rs,0,1,200;r1,1,0,1E3;j1,0,1,1E-5*v2;
            e1,0,3,100*v1;r2,3,2,1E4",2,0)
\end{entry}

Note that we have added only the resistance \code{rs}. The source
$V_s$ was not added because, to find the output impedance, the
voltage source must be dead, and a dead voltage source is simply a
short circuit. It is not necessary to place this short circuit in
series with the resistance, because placing only the resistance has
the same effect. We press \key{ENTER}. We select Passive and DC, and
press \key{ENTER}. We see that \code{zth} is 8571.

\problem{054}

\emph{Problem: Find the equivalent of the following network, in
transmission parameters.}

Solution: Below, the description we use:

\begin{entry}
sq\port("r1,1,3,2.;r2,3,2,4.;r3,3,0,10.",1,2)
\end{entry}

We press \key{ENTER}. We select \code{a}, some name and DC, and
press \key{ENTER}. We obtain the following transmission parameters:
1.2, 6.8, .1, 1.4.

\problem{055}

\emph{Problem: Find the equivalent of the following network, in
transmission parameters.}

Solution: Note that the network, in this case, is composed of two
groups of resistances in a T arrangement. Below, the description we
use:

\begin{entry}
sq\port("r1,1,2,2.;r2,2,0,10.;r3,2,3,4.;r4,3,4,4.;
        r5,4,0,20.;r6,4,5,8.",1,5)
\end{entry}

We press \key{ENTER}. We select \code{a}, some name and DC, and
press \key{ENTER}. We obtain the following transmission parameters:
1.78, 25.84, .19, 3.32. This is an example of how, with the
\code{port} tool, a network that would usually have taken two
two-ports can be reduced to a single two-port.

\problem{056}

\emph{Problem: Find the equivalent of the following network, in
transmission parameters.}

Solution: Below, the description we use:

\begin{entry}
sq\port("r1,1,3,5.;j1,3,0,.1*v2;e1,2,3,.3*v1;r2,2,0,4.",1,2)
\end{entry}

We press \key{ENTER}. We select \code{a}, some name and DC, and
press \key{ENTER}. We obtain the following transmission parameters:
2.12, 3.85, .35, 1.
